\textbf{1.} Sur une figure soignée, les trois médianes paraissent passer par
un même point.

\textbf{2.} Un tracé a une épaisseur et une précision limitées : trois droites
qui se coupent en trois points très proches (formant un minuscule triangle)
semblent concourantes. De plus, on ne trace qu'un nombre fini de triangles.

\textbf{3.} Soit $A'$, $B'$, $C'$ les milieux respectifs de $[BC]$, $[CA]$,
$[AB]$, et $G$ le point défini par
$\overrightarrow{GA} + \overrightarrow{GB} + \overrightarrow{GC} = \vec 0$
(il existe et il est unique : en introduisant $A$, cette égalité équivaut à
$\overrightarrow{AG} = \frac13\bigl(\overrightarrow{AB} + \overrightarrow{AC}\bigr)$).
Comme $A'$ est le milieu de $[BC]$,
$\overrightarrow{AB} + \overrightarrow{AC} = 2\,\overrightarrow{AA'}$, donc
\[
\overrightarrow{AG} = \frac23\,\overrightarrow{AA'} :
\]
$G$ appartient à la médiane $(AA')$. Le point $G$ étant défini de façon
symétrique en $A$, $B$, $C$, le même calcul montre que
$\overrightarrow{BG} = \frac23\,\overrightarrow{BB'}$ et
$\overrightarrow{CG} = \frac23\,\overrightarrow{CC'}$. Les trois médianes
passent donc par $G$, situé aux deux tiers de chacune à partir du sommet.
\hfill $\blacksquare$
